\documentclass[a4paper,11pt]{article}
\usepackage[utf8]{inputenc}
\usepackage[T1]{fontenc}
\usepackage[french]{babel}
\usepackage[margin=1.5cm]{geometry}
\usepackage{amsmath,amsfonts,amssymb}
\usepackage{tabularx}
\usepackage{enumitem}

% Ajustement de la hauteur des lignes du tableau pour laisser la place d'écrire
\renewcommand{\arraystretch}{2}

\begin{document}

\pagestyle{empty}

% ==========================================
% FICHE 2
% ==========================================
\begin{center}
    \Large\textbf{Entraînement : Automatismes 02}
\end{center}

\noindent \textbf{Partie 1 \hfill 6 points \hfill 20 minutes}

\medskip
\noindent \textbf{Consigne :} Pour chaque question, recopier sur la copie son numéro et la réponse correspondante. Aucune justification n'est demandée.

\bigskip

\noindent
\begin{tabularx}{\textwidth}{|c|X|p{4cm}|}
\hline
\textbf{N$^{\circ}$} & \textbf{Énoncé} & \textbf{Réponse} \\ \hline
\textbf{1} & Quel est le tiers de 48 ? & \\ \hline
\textbf{2} & Un film dure 135 min. Quelle est sa durée en heures et minutes ? & \\ \hline
\textbf{3} & Voici une série de notes : 8 ; 14 ; 11 ; 5 ; 12. Quelle est la médiane ? & \\ \hline
\textbf{4} & Dans un triangle $GHI$ rectangle en $H$, on sait que $\widehat{G} = 37^{\circ}$. Calculer la mesure de l'angle $\widehat{I}$. & \\ \hline
\textbf{5} & Sur une droite graduée, l'abscisse de $C$ est $\frac{5}{3}$ et celle de $D$ est $\frac{7}{6}$. Quel point est le plus proche de l'origine ? & \\ \hline
\textbf{6} & Dans le triangle $RST$ rectangle en $S$, quel est le sinus de l'angle $\widehat{SRT}$ ? \textit{(Donner la formule avec les côtés)} & \\ \hline
\textbf{7} & Un jeu vidéo coûte 40~€. Son prix augmente de 10~\%. Quel est le nouveau prix ? & \\ \hline
\textbf{8} & Dans un club de sport de 250 adhérents, 40~\% pratiquent la natation. Combien d'adhérents font de la natation ? & \\ \hline
\textbf{9} & Compléter le script Scratch pour dessiner un carré : \newline répéter \dots\dots fois \newline \hspace*{0.5cm} avancer de 60 \newline \hspace*{0.5cm} tourner de \dots\dots degrés & \\ \hline
\textbf{10} & Un segment $[KL]$ mesure 8 cm. On lui applique un agrandissement de rapport 1,5. Quelle est la longueur du segment agrandi $K'L'$ ? & \\ \hline
\end{tabularx}

\newpage

\subsection*{Correction (pour s'auto-évaluer)}
\begin{enumerate}[label=\textbf{\arabic*.}]
    \item \textbf{16} (car $48 \div 3 = 16$).
    \item \textbf{2 h 15 min} (car $135 = 2 \times 60 + 15$).
    \item \textbf{11} (Série ordonnée : 5 ; 8 ; \textbf{11} ; 12 ; 14).
    \item \textbf{53$^{\circ}$} (car $90 - 37 = 53$).
    \item \textbf{Le point D} (en mettant au même dénominateur, l'abscisse de $C$ est $\frac{10}{6}$, qui est supérieure à $\frac{7}{6}$).
    \item $\sin(\widehat{SRT}) = \mathbf{\frac{ST}{RT}}$ (Côté opposé / Hypoténuse).
    \item \textbf{44~€} (L'augmentation est de $40 \times 0,10 = 4$~€).
    \item \textbf{100 adhérents} (car $0,40 \times 250 = 100$).
    \item Répéter \textbf{4} fois et Tourner de \textbf{90} degrés.
    \item \textbf{12 cm} (car $8 \times 1,5 = 12$).
\end{enumerate}

\newpage

% ==========================================
% FICHE 3
% ==========================================
\begin{center}
    \Large\textbf{Entraînement : Automatismes 03}
\end{center}

\noindent \textbf{Partie 1 \hfill 6 points \hfill 20 minutes}

\medskip
\noindent \textbf{Consigne :} Pour chaque question, recopier sur la copie son numéro et la réponse correspondante. Aucune justification n'est demandée.

\bigskip

\noindent
\begin{tabularx}{\textwidth}{|c|X|p{4cm}|}
\hline
\textbf{N$^{\circ}$} & \textbf{Énoncé} & \textbf{Réponse} \\ \hline
\textbf{1} & Quel est le triple de 27 ? & \\ \hline
\textbf{2} & Un train part à 8 h 45 et arrive à 11 h 10. Quelle est la durée totale du trajet ? & \\ \hline
\textbf{3} & Quelle est l'étendue de la série de valeurs suivante : 6 ; 18 ; 12 ; 9 ; 15 ? & \\ \hline
\textbf{4} & Dans un triangle $MNP$, on sait que $\widehat{M} = 50^{\circ}$ et $\widehat{N} = 65^{\circ}$. Quelle est la nature précise du triangle $MNP$ ? & \\ \hline
\textbf{5} & Donner le résultat sous la forme d'une fraction irréductible : $\frac{3}{5} + \frac{1}{10}$. & \\ \hline
\textbf{6} & Dans le triangle $IJK$ rectangle en $J$, quelle est la tangente de l'angle $\widehat{JIK}$ ? \textit{(Donner la formule avec les côtés)} & \\ \hline
\textbf{7} & Un pull affiché à 50~€ est soldé avec une remise de 30~\%. Quel est son prix soldé ? & \\ \hline
\textbf{8} & Sur un parking de 60 véhicules, 15 sont électriques. Quel pourcentage du nombre total de véhicules cela représente-t-il ? & \\ \hline
\textbf{9} & Compléter le script Scratch pour dessiner un hexagone régulier : \newline répéter \dots\dots fois \newline \hspace*{0.5cm} avancer de 40 \newline \hspace*{0.5cm} tourner de \dots\dots degrés & \\ \hline
\textbf{10} & L'aire d'un triangle est de $10\text{ cm}^2$. On lui applique un agrandissement de rapport $k = 3$. Quelle est l'aire du triangle obtenu ? & \\ \hline
\end{tabularx}

\newpage

\subsection*{Correction (pour s'auto-évaluer)}
\begin{enumerate}[label=\textbf{\arabic*.}]
    \item \textbf{81} (car $27 \times 3 = 81$).
    \item \textbf{2 h 25 min} (de 8h45 à 9h : 15 min ; de 9h à 11h10 : 2 h 10 min).
    \item \textbf{12} (Valeur maximale $-$ Valeur minimale $= 18 - 6 = 12$).
    \item \textbf{Triangle isocèle en P} (L'angle $\widehat{P}$ mesure $180 - (50 + 65) = 65^{\circ}$, il y a donc deux angles de $65^{\circ}$).
    \item $\mathbf{\frac{7}{10}}$ (car $\frac{3}{5} = \frac{6}{10}$ et $\frac{6}{10} + \frac{1}{10} = \frac{7}{10}$).
    \item $\tan(\widehat{JIK}) = \mathbf{\frac{JK}{IJ}}$ (Côté opposé / Côté adjacent).
    \item \textbf{35~€} (La remise est de $50 \times 0,30 = 15$~€, d'où $50 - 15 = 35$).
    \item \textbf{25~\%} (La proportion est $\frac{15}{60} = \frac{1}{4} = 0,25$).
    \item Répéter \textbf{6} fois et Tourner de \textbf{60} degrés (car $360 \div 6 = 60$).
    \item $\mathbf{90\text{ cm}^2}$ (Dans un agrandissement de rapport $k$, l'aire est multipliée par $k^2$, soit ici par $3^2 = 9$).
\end{enumerate}

\end{document}